\section{As the Sun Brings Fire}

\textbf{Petrarch} wrote a series of sonnets to his twin flame Laura, although their real meaning is esoteric. The sonnet has this structure:

\begin{enumerate}
\item The first quatrain indicates a conflict or problem

\item The second quatrain explains the problem

\item The next three lines, the ``volta'', changes the attitude or perspective

\item The final three lines are the resolution
\end{enumerate}

The initiate into the Fedeli d'Amore would be challenged with a sonnet, and he would have to answer back with an appropriate sonnet of his own. It is worth the trouble to learn the art, because you never know.

I pondered \textit{Sonnet 164: Heaven and Earth are Silent}\footnote{\url{https://gornahoor.net/?p=15071}}. Who can still the heavens, earth, and the wind, subdue the animals, and calm the sea? Yet, internally he is at war with himself. The same hand brings him both distress and peace, pain and healing. With the wind calm, he is stuck at sea in the cycle of birth and death, sweet and bitter.

Who has power over the Sky, Earth, and Wind? Who brings sleep to the animals? Who rides her Chariot into the Night? It is Nyx, the primeval goddess of the night. Her alluring beauty and power which arises out of Chaos or Non-Being both attracts and frightens. Even Zeus feared her power.

A man's power is visible and out in the open. The power of Nyx, however, is hidden in the darkness; it is meonic. But where else can freedom, creativity, or love arise from? They are surprises because they have no physical or causal antecedent. Not all surprises are nice, as you know. Nyx has several children who are relentless and will never stay out of your life. Hence, playing it safe cannot prevent accidents, but will stifle freedom, creativity, and love.

\textbf{Nicholas Berdyaev} recognizes the power of Nyx without naming her:

\begin{quotex}
uncreated freedom, non-being which is prior to being, the meonic abyss which is neither Creator nor creature ... This is the ultimate mystery behind reality. Endless consequences follow from it. It accounts both for evil and for the creation of what has never existed before. The ethics of creativeness goes back to this primary truth.
\end{quotex}


Love also comes out of the night, hidden by darkness, so it is a surprise when it suddenly appears. It lights up knowledge. Then there is the knower and the known. Esoterically, knowing is the same as being, so in cosmic love, the knower and the known are One being, entangled throughout time and space. This is God's will as Berdyaev explains:

\begin{quotex}
God's conception of man is a complete, masculinely feminine being, solar and tellurgic, logoic and cosmic at the same time. Only insofar as he is complete is he chaste, wise, and Sophian in his perfect wholeness. As a sexual halved, divided being he is not chaste, or wise, and is doomed to disharmony, to passionate longing and dissatisfaction.
\end{quotex}


The following sonnet to Nyx was dictated to me in seven syllables rather than iambic pentameter. Seven syllables are an option in Italian sonnets, so I've decided to keep it this way.

This meditation was accompanied by some two strong dreams this week. The first was an encounter with the Puer in a quite vivid, even disconcerting, way. I fell back asleep was awakened by the very same dream 40 minutes later. I don't recall ever having a dream of the Puer previously.

The next night I dreamed of the Shadow of the Puer. He was hiding in the dark garage. I knew but entered anyway. He threw something at me as if to get my attention and scampered away into the unknown. I woke up screaming and could not fall back asleep for hours. Now, today, I don't know if I was the Puer in the first dream or the Shadow of the Puer. This is what \textbf{Carl Jung} writes about the encounter with the Puer in \emph{Answer to Job}.

\begin{quotex}
That higher and `complete' man is begotten by the `unknown' father and born from Wisdom, and it is he who, in the figure of the \textbf{puer aeternus} --- \emph{of changeful countenance, both white and black} --- represents our totality, which transcends consciousness. It was this boy into whom Faust had to change, abandoning his inflated onesidedness which saw the devil only outside. Christ's `Except ye become as little children' prefigures this change, for in them the opposites lie close together; but what is meant is the boy who is born from the maturity of the adult man, and not the unconscious child we would like to remain.
\end{quotex}


\paragraph{Nyx, the goddess of the night}

\begin{verse}
Nyx, the goddess of the night,\\
Secret source of all power\\
Zeus himself had feared her might\\
Her beauty tempts the hour.

When sane men shun the abyss\\
Where men's souls are judged aright.\\
A glimpse before Sleep brings bliss\\
Or else Death's cold breath brings fright.

Ten thousand times appeared she,\\
My life faded in her dream\\
Geras cannot quell desire.

Distress and strife torment me\\
Prophecies aren't what they seem,\\
Night fades as the Sun brings fire.
\end{verse}

Some background on Nyx if you are unfamiliar with her.

\url{https://www.youtube.com/watch?v=Pw1OgUZi0tg}

\flrightit{Posted on 2021-10-17 by Cologero}

\begin{center}* * *\end{center}

\begin{footnotesize}
\begin{sffamily}
\texttt{Paulo Adolpho on 2021-10-19 at 23:55 said:}

Cologero, I'm tired of seeing the weakness of the masculine towards the feminine, in this case, according to the text, even Zeus is powerless and fearful. To you, is there a solution to all this hell, or are we destined to be powerless and slaves to that sex?

\hfill

\texttt{Cologero on 2021-10-20 at 21:26 said:}

Only the weak are so destined.

\hfill

\texttt{Paulo Adolpho on 2021-10-20 at 22:45 said:}

What would make a man powerful enough, to make him deal with a woman in a impeccable way?


\hfill

\end{sffamily}
\end{footnotesize}

